\answer{ex:02:1}{Hàng dài $3.2$~mm với bước $25\um$: $\approx129$ bộ phát hiện. Phát xung là một dải rộng $v\,\tau\ln2=1.7$~mm --- $\approx69$ bộ phát hiện, hơn nửa hàng; hướng là tâm của dải.}
\answer{ex:02:2}{Cửa sổ $-\tau\ln\frac{w_2}{\theta-w_1}\le\delta\le\tau\ln\frac{w_1}{\theta-w_2}$, tức là $[-0.235,\,0.178]\ms$; tâm $c=\frac\tau2\ln\frac{w_1(\theta-w_1)}{w_2(\theta-w_2)}=-28\,\mu\text{s}$. Hướng lệch $28\,\mu\text{s}$ về phía tai mạnh hơn --- gần ba bước phân giải.}
\answer{ex:02:3}{Cần các $s_i$ với $s_iw_{ij}s_j\ge0$; chúng tồn tại khi và chỉ khi mọi chu trình đều chẵn. Ức chế lẫn nhau quy được (không có dao động bền vững), vòng \nt{GLUT}--\nt{GABA} thì không (có thể dao động).}
\answer{ex:02:4}{Sáu nơ-ron $g_k=[x+y+z\ge k]$ và $\bar g_k=[\bar x+\bar y+\bar z\ge4-k]$, $k=1,2,3$; $C=g_2$, $\bar C=\bar g_2$, $S=[g_1+\bar g_2+g_3\ge2]$, $\bar S=[\bar g_1+g_2+\bar g_3\ge2]$: tám nơ-ron. Bằng AND và OR --- $12$.}
\answer{ex:02:5}{$T_{\min}(0.6)=0.718\,\tau$; $I_c=0.225$, gần ngưỡng $T_{\min}\propto(I-I_c)^{-1/2}$.}
\answer{ex:02:6}{(a)~$500$ mắt xích, $5\cdot10^4$ nơ-ron; độ phân tán $4.5\ms$ ($0.45\,\%$), sai số tương đối $\propto T^{-1/2}$. (b)~Một hệ số ngẫu nhiên của độ trễ, chung cho mọi mắt xích, với độ phân tán $5\,\%$.}
\answer{ex:02:7}{Làm tươi mỗi $\tau_F\ln2=0.69\,\text{s}$ một lần: $4.3$ xung mỗi giây trên một nơ-ron so với $20$, tiết kiệm hơn $4.6$ lần ($4\cdot10^{-7}$ so với $2\cdot10^{-6}$~J với $10^{-10}$~J mỗi xung). Flip-flop mất bit; tụ điện giữ được bit nếu việc làm câm bắt đầu không muộn hơn $0.49\,\text{s}$ sau lần làm tươi (xác suất $0.71$).}
